\subsection{Positive and negative parts}

DEFINITION 4. --- In a lattice ordered group G the positive part (resp. negative part, absolute value) of an element \(x \in G\) is the element \(\sup(x, 0)\) (resp. \(\sup(-x, 0)\) , \(\sup(x, -x)\) ), which is denoted \(x^{+}\) (resp. \(x^{-}\) , \(|x|\) ).

\pandocbounded{\includegraphics[keepaspectratio]{images/part002_7ee71349.jpg}}

Despite its name, the negative part \(x^{-}\) is a positive element.

Clearly \(x^{-} = (-x)^{+}\) and \(|-x| = |x|\) . Let us also note the following formulae, the first of which is an immediate consequence of the definitions and of the invariance of the ordering under translation, and the second of which follows from the first by Prop. 7 of VI, p.~10:

\[
\left\{ \begin{array}{l} \sup (x, y) = x + (y - x) ^ {+}, \\ \inf (x, y) = y - (y - x) ^ {+}. \end{array} \right.\tag{3}
\]

PROPOSITION 9. --- a) For each element x of a lattice ordered group G we have \(x = x^{+} - x^{-}\) and \(\inf(x^{+}, x^{-}) = 0\) .

\begin{enumerate}
\def\labelenumi{\alph{enumi})}
\setcounter{enumi}{1}
\item
  For every expression of \(x\) as the difference of two positive elements, say \(x = u - v\) , we have \(u = x^{+} + w\) and \(v = x^{-} + w\) with \(w = \inf(u, v)\) . In particular if \(\inf(u, v) = 0\) then \(u = x^{+}\) and \(v = x^{-}\) .
\item
  The relation \(\ll x\leqslant y\gg\) is equivalent to \(\ll x^{+}\leqslant y^{+}\) and \(x^{-}\geqslant y^{-}\gg\) .
\item
  We have \(|x| = x^{+} + x^{-} \geqslant 0\) .
\item
  For any \(x\) and \(y\) in \(G\) , we have the inequality \(|x + y| \leqslant |x| + |y|\) , and more generally \(\left|\sum_{i=1}^{n} x_i\right| \leqslant \sum_{i=1}^{n} |x_i|\) for any finite family \((x_i)\) of elements of \(G\) .
\item
  For any \(x\) and \(y\) in \(G\) we have \(||x| - |y|| \leqslant |x - y|\) .
\end{enumerate}

We will prove \(a)\) and \(b)\) simultaneously. If \(x = u - v\) with \(u\) and \(v\) positive, then \(u \geqslant x\) , so \(u \geqslant \sup (x, 0) = x^{+}\) , and \(w = u - x^{+}\) is positive. On the other hand we have

\[
x ^ {+} - x = \sup (x, 0) - x = \sup (x - x, - x) = x ^ {-}
\]

from which it follows that \(x = x^{+} - x^{-}\) , and \(v - x^{-} = w\) . If \(z \leqslant x^{-}\) then \(z \leqslant x^{+} - x\) , and so \(x \leqslant x^{+} - z\) ; if also \(z \leqslant x^{+}\) , then \(x^{+} - z\) is positive, and so \(x^{+} \leqslant x^{+} - z\) by definition of \(x^{+}\) . Hence we have \(z \leqslant 0\) , which implies \(\inf(x^{+}, x^{-}) = 0\) , whence by translation \(\inf(u, v) = w\) .

\begin{enumerate}
\def\labelenumi{\alph{enumi})}
\setcounter{enumi}{2}
\item
  The relation \(x \leqslant y\) implies \(\sup (y, 0) \geqslant x\) and \(\sup (y, 0) \geqslant 0\) , hence \(x^{+} \leqslant y^{+}\) ; similarly if \(-y \leqslant -x\) we deduce \(x^{-} \geqslant y^{-}\) . The converse implication follows immediately from \(x = x^{+} - x^{-}\) and \(y = y^{+} - y^{-}\) .
\item
  Since \(x \leqslant x^{+}\) and \(-x \leqslant x^{-}\) , it is clear that
\end{enumerate}

\[
\left| x \right| = \sup (x, - x) \leqslant x ^ {+} + x ^ {-}.
\]

Conversely, if \(a \geqslant x\) and \(a \geqslant -x\) then we deduce from \(c\) ) that \(a^+ \geqslant x^+\) , \(a^+ \geqslant x^-\) , \(a^- \leqslant x^-\) and \(a^- \leqslant x^+\) ; since \(a^-\) is positive and \(\inf(x^+, x^-) = 0\) , the last two inequalities imply \(a^- = 0\) and \(a = a^+\) ; the first two inequalities then give \(a \geqslant \sup(x^+, x^-)\) , and \(\sup(x^+, x^-)\) is equal to \(x^+ + x^-\) by \(a\) ) and by Prop. 7 of VI, p.~10.

\begin{enumerate}
\def\labelenumi{\alph{enumi})}
\setcounter{enumi}{4}
\item
  Since \(x \leqslant |x|\) and \(y \leqslant |y|\) , we have \(x + y \leqslant |x| + |y|\) ; since \(-x \leqslant |x|\) and \(-y \leqslant |y|\) , we have \(-x - y \leqslant |x| + |y|\) ; whence the first inequality. The second follows by induction on \(n\) .
\item
  Replacing \(x\) and \(y\) by \(y\) and \(x - y\) in \(e\) ), we obtain
\end{enumerate}

\[
\left| x \right| - \left| y \right| \leqslant \left| x - y \right|;
\]

similarly we have \(|y| - |x| \leqslant |y - x| = |x - y|\) ; whence the stated result.

Remark. --- We deduce from \(d\) that \(|x|=0\) implies \(x=0\) (for \(x^{+}\) and \(x^{-}\) are positive); thus \(x\neq0\) implies \(|x|>0\) .

PROPOSITION 9 (DIV). --- If the group \(\mathcal{P}^*\) of principal fractional ideals of K is lattice ordered, then every element \(x\) of \(\mathbf{K}^*\) can be written in the form \(x = uv^{-1}\) , where \(u\) and \(v\) are elements of A such that \(1 = \gcd(u, v)\) ; for any other expression \(x = u'v'^{-1}\) of \(x\) as the quotient of two elements of A, we have \(u' = uw\) and \(v' = vw\) , where \(w\) is a gcd of \(u'\) , \(v'\) ; in particular if \(1 = \gcd(u', v')\) then \(u'\) and \(v'\) are associates of \(u\) and \(v\) respectively.

Such an expression \(uv^{-1}\) for an element x of \(K^{*}\) is often called a reduced fraction.

